\scanpage{18}
% archive sheet number: 17

% slanted
\margin{\deleted{\unsure{c'est ici que|c'est} le fait que les $H_i$ soient distincts est \unsure{essentiel|inessentiel}}}

On trouve que c'est l'intersection des $\tfrac12$ espaces \uline{fermés} $\overline{V'_{H_i}}=V'_{H_i}\cup H_i$, $\overline{V''_{H_i}}=V''_{H_i}\cup H_i$, déterminés par les $H_i$ ($i\in I$), \inserted{qui contiennent \unsure{$Z_i$|$Z$},} en d'autres termes

\[
\overline{Z^*_\alpha}=\bigcup_{\substack{\beta\in\Sigma\\\text{tel que }i\in I\\\text{implique }k_i(\beta)\leq k_i(\alpha)\\\deleted{\unsure{\sim}}}}Z^*_\beta
=\bigcup_{\substack{\beta\in\Sigma\\\beta\leq\alpha}}Z^*_\beta
\]

où les \deleted{\unsure{$k_i$}}

\[
k_i:E\longrightarrow\deleted{E}D_{H_i}\qquad(i\in I)
\]

sont les fonctions coordonnées, et où on introduit dans $\deleted{E}D_{H_i}$ la relation d'ordre $\leq$ donnée par

\diagram{\begin{tikzcd}[ampersand replacement=\&,column sep=tiny,row sep=small]
\unsure{a'_i|a'|\alpha'_i}\arrow[dr]\&\&\unsure{a''_i|a''|\alpha''_i}\arrow[dl]\\
\&0\&
\end{tikzcd}}

$a',a''$ les deux \unsure{élémts|éléments} $\ne 0$ de $D_{\unsure{H_i|H}}$

i.e. Si $x,y\in D_{H_i}$ on a $x\leq y$ sss $x=0$ ou $x=y$, et on met sur $E=\prod_{i\in I}D_{\unsure{I|i}}$ l'ordre produit.

On \deleted{est alors} met sur $\Sigma$ l'ordre induit, d'où une ``top. inférieure'' sur $\Sigma$, dont les fermés sont les parties \inserted{$A$} de $\Sigma$ telles que $\alpha\in A,\ \beta\leq\alpha\Rightarrow\beta\in A$. On a donc

\[
\beta\leq\alpha\iff\beta\in\overline{\{\alpha\}},\qquad\text{et}
\]

\[
\overline{Z^*_\alpha}=Z^*_{\overline{\{\alpha\}}}
\]

en posant, pour $A\subset\Sigma$

\[
Z^*_A=\bigcup_{\alpha\in A}Z^*_\alpha\ .
\]
